\origpage{330}{N. Robertson, P.D. Seymour / Journal of Combinatorial Theory, Series B 92 (2004) 325--357}
\noindent patchworks. Then $C$ is also a set of partial $\Omega_2$-patchworks; and it is an easy exercise to show that $C$ is well-behaved taking $\Omega = \Omega_1$, if and only if it is well-behaved with $\Omega = \Omega_2$. Thus, our terminology suppressing the dependence on $\Omega$ is not misleading.

The following is Theorem 6.7 of [4].

\textbf{3.1.} \textit{Let $\Omega$ be a well-quasi-order, let $\mathcal{F}$ be a well-behaved set of rootless partial $\Omega$-patchworks, and let $\theta \ge 1$ be an integer. Let $P_i = (G_i, \mu_i, \Delta_i, \phi_i)$ ($i = 1, 2, \dots$) be a countable sequence of rootless robust $\Omega$-patchworks. For each $i \ge 1$ let $\lambda_i$ be a tie-breaker in $G_i$; and suppose that for every tangle $\mathcal{T}$ in $G_i$ of order $\ge \theta$ there is a rooted location $\mathcal{L}$ in $G_i$ such that $\mathcal{L}$ $\theta$-isolates $\mathcal{T}$ and $(P_i, \mathcal{L})$ has a heart in $\mathcal{F}$. Then there exist $j > i \ge 1$ such that $P_i$ is simulated in $P_j$.}

\sectionhead{4. Well-behaved sets of patchworks}

The previous result 3.1, combined with the main result of [3] (see 10.3 of the present paper), almost proves Wagner's conjecture. Not quite, however; although the rooted locations provided by [3] have hearts in a well-behaved set, they do not quite $\theta$-isolate the corresponding tangles and so 3.1 cannot be applied to them. In the next few sections we prove a strengthening 7.3 of 3.1, that bridges the gap. We show that the locations of [3] can be modified such that the new locations still have hearts in a (new) well-behaved set and do $\theta'$-isolate the corresponding tangles, for an appropriate $\theta'$. The main problem is that there are a bounded number of vertices that need to be removed; and in essence 7.3 addresses the problems caused by removing these vertices.

To prove 7.3, we first need to develop ways of constructing new well-behaved sets of patchworks from old ones, and that is the object of this section. Incidentally, the rooted locations $\mathcal{L}$ provided by [3] have the property that $\bigcup(A^- : A \in \mathcal{L}) = G^-$, which has two desirable consequences; that their hearts have no ``isolated vertices'', and that their hearts have no edges labeled from $\Omega$, and hence are more naturally regarded as patchworks than as partial $\Omega$-patchworks. This motivates the following.

If $P = (G, \mu, \Delta)$ is a patchwork and $\Omega$ is a quasi-order, we call every $\Omega$-patchwork $(G, \mu, \Delta, \phi)$ an $\Omega$\textit{-completion} of $P$. A set $\mathcal{F}$ of patchworks is \textit{well-behaved} if for every well-quasi-order $\Omega$ and every countable sequence $P_i$ ($i = 1, 2, \dots$) of $\Omega$-completions of members of $\mathcal{F}$ there exist $j > i \ge 1$ such that $P_i$ is simulated in $P_j$.

\textbf{4.1.} \textit{If $\mathcal{F}$ is well-behaved, then there exists $N \ge 0$ such that if $(G, \mu, \Delta) \in \mathcal{F}$ and $e \in dom(\mu)$ then $|\bar{\mu}(e)| \le N$.}

\textbf{Proof.} Let $\Omega$ be the well-quasi-order with $E(\Omega) = \{\omega_1, \omega_2\}$ say, where $\omega_1, \omega_2$ are incomparable (that is, $\omega_1 \not\le \omega_2 \not\le \omega_1$). Suppose that there is no $N$ as in the theorem. Then there exist integers $n_i$ and $P_i = (G_i, \mu_i, \Delta_i) \in \mathcal{F}$ and $e_i \in E(G_i) \cap dom(\mu_1)$ with $|\bar{\mu}_i(e_i)| = n_i$ for $i = 1, 2, \dots$, such that $n_1 < n_2 < \dots$. For $i \ge 1$, define $\phi_i : E(G_i) \to E(\Omega)$ by $\phi_i(e_i) = \omega_2$ and $\phi_i(e) = \omega_1$ ($e \ne e_i$). Then $(G_i, \mu_i, \Delta_i, \phi_i)$ ($= Q_i$, say) is an $\Omega$-completion of $P_i$. Since $\mathcal{F}$ is well-behaved, there exist $j > i \ge 1$ such that there is a realizable expansion $\eta$ of
